%% bare_jrnl.tex
%% V1.4b
%% 2015/08/26
%% by Michael Shell
%% see http://www.michaelshell.org/
%% for current contact information.
%%
%% This is a skeleton file demonstrating the use of IEEEtran.cls
%% (requires IEEEtran.cls version 1.8b or later) with an IEEE
%% journal paper.
%%
%% Support sites:
%% http://www.michaelshell.org/tex/ieeetran/
%% http://www.ctan.org/pkg/ieeetran
%% and
%% http://www.ieee.org/
%%
%% ---------------------------------------------------------------
%% Modified by Abstrack (TFG).
%% Plantilla de referencia: IEEE Journal LaTeX Template oficial
%% (IEEEtran.cls, distribuido vía template-selector.ieee.org),
%% con el front-matter ajustado a las Author Guidelines de
%% IEEE Internet of Things Journal (ieee-iotj.org/guidelines-for-authors).
%% Abstract e introducción se insertan automáticamente.
%% ---------------------------------------------------------------

%%*************************************************************************
%% Legal Notice:
%% This code is offered as-is without any warranty either expressed or
%% implied; without even the implied warranty of MERCHANTABILITY or
%% FITNESS FOR A PARTICULAR PURPOSE!
%% User assumes all risk.
%% In no event shall the IEEE or any contributor to this code be liable for
%% any damages or losses, including, but not limited to, incidental,
%% consequential, or any other damages, resulting from the use or misuse
%% of any information contained here.
%%
%% All comments are the opinions of their respective authors and are not
%% necessarily endorsed by the IEEE.
%%
%% This work is distributed under the LaTeX Project Public License (LPPL)
%% ( http://www.latex-project.org/ ) version 1.3, and may be freely used,
%% distributed and modified. A copy of the LPPL, version 1.3, is included
%% in the base LaTeX documentation of all distributions of LaTeX released
%% 2003/12/01 or later.
%% Retain all contribution notices and credits.
%% ** Modified files should be clearly indicated as such, including  **
%% ** renaming them and changing author support contact information. **
%%*************************************************************************


\documentclass[journal]{IEEEtran}

\hyphenation{op-tical net-works semi-conduc-tor}


\begin{document}

\title{[ARTICLE TITLE]}

\author{[FIRST AUTHOR],~\IEEEmembership{Member,~IEEE,}
        and~[CO-AUTHOR]%
\thanks{[AFFILIATION DATA].}%
\thanks{Manuscript received [DATE].}}

\markboth{IEEE INTERNET OF THINGS JOURNAL,~Vol.~[VOL], No.~[NUM], [MONTH]~[YEAR]}%
{[AUTHORS]: [SHORT TITLE]}

\maketitle

% IEEE IoT-J: el abstract debe tener entre 150 y 250 palabras y no debe
% incluir citas bibliográficas ni ecuaciones.
\begin{abstract}
Este trabajo aborda las limitaciones de modelos recurrentes, como RNN, LSTM y GRU, que presentan dependencias secuenciales que restringen la paralelización y reducen la eficiencia computacional, dificultando modelar dependencias a largo plazo en tareas como traducción automática. Se propone el Transformer, arquitectura que elimina la recurrencia y usa solo mecanismos de atención para capturar dependencias globales entre entradas y salidas. El Transformer usa una estructura encoder-decoder con seis capas idénticas en el modelo base, combinando atención multi-cabeza y redes feed-forward aplicadas por posición, junto con codificaciones posicionales sinusoidales que incorporan el orden sin recurrencia ni convoluciones. Se entrena con Adam y técnicas de regularización como dropout y label smoothing, usando datos masivos como el WMT 2014 para traducción inglés-alemán e inglés-francés, en hardware de 8 GPUs NVIDIA P100 durante 12 horas a 3.5 días según el tamaño. Los resultados muestran que el Transformer grande alcanza un BLEU de 28.4 en inglés-alemán, superando en más de 2 puntos a los mejores modelos previos y ensambles, y un BLEU de 41.0 en inglés-francés, mejorando modelos individuales anteriores con menor costo de entrenamiento. El modelo base también supera a competidores anteriores a menor costo. En parsing sintáctico en inglés, alcanza un F1 de hasta 92.7 en configuración semi-supervisada, superando métodos recurrentes clásicos. Experimentos con variaciones evidencian el impacto del número de cabezas y dimensiones internas en el rendimiento. En conclusión, el Transformer es el primer modelo de transducción que prescinde totalmente de recurrencia, permitiendo un entrenamiento más rápido y resultados superiores en traducción y análisis sintáctico, con perspectivas de extensión a otras modalidades y mecanismos de atención local para secuencias mayores.
\end{abstract}

\begin{IEEEkeywords}
[KEYWORD 1], [KEYWORD 2], [KEYWORD 3]
\end{IEEEkeywords}

\IEEEpeerreviewmaketitle


\section{Introduction}

En el campo de los modelos de transducción de secuencias, las arquitecturas dominantes tradicionalmente han utilizado redes neuronales recurrentes (RNN) o convolucionales (CNN) que incluyen mecanismos de atención para conectar el codificador y el decodificador. Estos modelos han logrado avances significativos en tareas como la traducción automática, donde la eficiencia computacional y la calidad de la traducción son aspectos fundamentales. La atención se ha consolidado como un componente clave para mejorar el rendimiento al permitir que el modelo se enfoque dinámicamente en diferentes partes de la secuencia de entrada durante la generación de la salida.

Dentro de este panorama, la búsqueda por mejorar la eficiencia y la capacidad de paralelización de los modelos ha sido un tema central. Las arquitecturas basadas en recurrencias o convoluciones presentan limitaciones inherentes, como la necesidad de cómputos secuenciales o la dificultad para capturar dependencias a largo plazo de manera efectiva. En este contexto, surge la relevancia de explorar nuevas aproximaciones que permitan superar estos obstáculos, manteniendo o mejorando la calidad en las tareas de procesamiento de lenguaje natural, específicamente en traducción automática. El Transformer se posiciona como una innovadora alternativa que prescinde por completo de recurrencias y convoluciones, utilizando exclusivamente mecanismos de atención para tratar las secuencias.

A pesar de los avances en modelos basados en atención para la traducción neuronal, persisten limitaciones importantes. Bahdanau et al. (2015) [VERIFICAR] presentan un modelo que mejora la capacidad de los modelos encoder-decoder tradicionales al evitar comprimir toda la oración de entrada en un vector fijo, lo que dificulta el manejo de oraciones largas. Sin embargo, esta arquitectura sigue usando redes recurrentes, lo que limita la eficiencia computacional y la posibilidad de paralelización efectiva.

Por otro lado, Sutskever et al. (2014)~\cite{sutskever2014seq2seq} desarrollan un modelo sequence-to-sequence basado en LSTM profundos que logra mejoras en la traducción, pero depende de representar la secuencia completa en vectores de dimensión fija y requiere invertir la secuencia de entrada para manejar mejor las dependencias largas. Esta estrategia, aunque efectiva, implica desafíos en la optimización y procesamiento.

En conjunto, estas limitaciones señalan la necesidad de una arquitectura que supere las restricciones propias de las redes recurrentes y convolucionales, principalmente la secuencialidad en el procesamiento y la dificultad para paralelizar, y que facilite un entrenamiento más eficiente y un manejo adecuado de las dependencias a largo plazo sin requerir codificaciones fijas. Este es el vacío que justifica la exploración de enfoques alternativos como el Transformer.

La propuesta central del trabajo es el Transformer, una arquitectura que responde directamente a la necesidad de superar las limitaciones de eficiencia y paralelización presentes en las arquitecturas basadas en redes recurrentes y convolucionales. Este modelo elimina completamente el uso de recurrencias y convoluciones, sustituyéndolos por un mecanismo de auto-atención multi-cabeza que permite establecer conexiones directas entre cualquier par de posiciones dentro de una secuencia.

Este enfoque es especialmente adecuado para abordar los problemas asociados a la secuencialidad en el cómputo y la dificultad de paralelización, ya que reduce el coste computacional a una cantidad constante en función de la profundidad del modelo. De esta manera, se posibilita un entrenamiento más eficiente, con una mejor capacidad para capturar dependencias de largo alcance en las secuencias, sin la necesidad de codificaciones fijas que suelen limitar otros métodos. Así, el Transformer ofrece una solución innovadora que mejora el rendimiento en tareas de traducción automática, cumpliendo con las demandas identificadas en investigaciones previas.

- Introducción de la arquitectura Transformer basada completamente en atención multi-cabeza, que permite manejar las secuencias sin recurrir a recurrencias ni convoluciones (sección 3).
- Propuesta del mecanismo de atención escalar producto punto y multi-cabeza, diseñado para mejorar la representación y el procesamiento paralelo de las secuencias.
- Uso de codificaciones posicionales sinusoidales para incorporar información sobre el orden de los elementos en la secuencia sin depender de estructuras recurrentes (sección 3.5).
- Demostración experimental de la superioridad del modelo Transformer en tareas de traducción automática, específicamente en los conjuntos de datos WMT 2014 inglés-alemán e inglés-francés (sección 6.1).
- Aplicación del modelo para análisis de constituyentes en inglés, evidenciando la capacidad de generalización del enfoque más allá de la traducción (sección 6.3).
- Análisis y visualización de los mecanismos de atención, mostrando comportamientos interpretables que se relacionan con la estructura sintáctica y semántica del idioma (apéndice y sección 4).
- Publicación del código asociado al Transformer para facilitar la reproducibilidad y permitir que la comunidad científica utilice y extienda el modelo.

La validación del Transformer se realiza mediante experimentos en tareas estándar de traducción automática, específicamente en los conjuntos de datos WMT 2014 para traducción de inglés a alemán e inglés a francés, donde se compara su desempeño con modelos existentes y variantes del propio modelo para evaluar el impacto de diferentes hiperparámetros. Además, se extiende la evaluación a tareas de análisis sintáctico en inglés, demostrando la capacidad de generalización del modelo. Complementariamente, se llevan a cabo análisis e interpretaciones de los mecanismos de atención internos para comprender su funcionamiento y características explicativas, todo ello respaldando y justificando la solidez y eficacia de las contribuciones presentadas.

El resto del documento se organiza comenzando con la introducción y motivación del trabajo, seguida por una descripción detallada de la arquitectura del Transformer, una justificación y análisis de las ventajas del mecanismo de auto-atención, y la presentación del régimen de entrenamiento. A continuación, se exponen los resultados experimentales que evalúan el desempeño del modelo en diferentes tareas, para finalmente concluir con una discusión y planteamiento de líneas futuras de trabajo. Además, el documento incluye referencias y apéndices que complementan la información técnica y experimental presentada.


\section{[SECTION 2 TITLE]}
\label{sec:section2}

[Section 2 content.]


\section{[SECTION 3 TITLE]}
\label{sec:section3}

[Section 3 content.]


\section{[SECTION 4 TITLE]}
\label{sec:section4}

[Section 4 content.]


\section{Conclusion}
\label{sec:conclusion}

[Conclusions.]


\appendices
\section{Appendix}
[Appendix content.]


\section*{Acknowledgment}

[Acknowledgments.]


% Sin bibliografia: no se proporcionaron citas.


\begin{IEEEbiographynophoto}{[FIRST AUTHOR]}
[Brief author biography.]
\end{IEEEbiographynophoto}


\bibliographystyle{IEEEtran}
\bibliography{attention_is_all_you_need}

\end{document}
